\documentclass{article}
% \usepackage{amsmath}
\usepackage[fleqn]{amsmath}
\usepackage{amssymb}
% \usepackage{mathtools}
\usepackage{bussproofs}
\usepackage{mathrsfs}
\usepackage{svg}
\usepackage[margin=1in]{geometry}
\usepackage{stackengine}
\usepackage{titlesec}
\usepackage{fitch}
\usepackage{multicol}
\usepackage[bottom]{footmisc}
\usepackage[utf8x]{inputenc}
\usepackage{url}
\usepackage{hyperref}
\usepackage{listings}
\usepackage{inconsolata}


\lstset{
  basicstyle=\ttfamily\fontfamily{zi4}\selectfont,
  breaklines=true,
  escapeinside={/*!}{!*/},
}

\begin{document}

\section{TypeScript Grammar (WIP)}

\begin{lstlisting}
Type ::= AnyType
       | UnknownType
       | NeverType
       | VoidType
       | ObjectType
       | StringType
       | NumberType
       | BooleanType
       | SymbolType
       | BigIntType
       | ObjectLiteralType
       | Id
       | "(" Type ")" [bracket]
       > left:
         Type "&" Type
       > Type "|" Type

ObjectLiteralType ::= "{" PropertySignatures "}"
PropertySignature ::= Id ":" Type
                    | Id "?" ":" Type
                    | "readonly" Id ":" Type
                    | "readonly" Id "?" ":" Type

Identifier ::= see https://tc39.es/ecma262/#prod-IdentifierName

TypeAliasDeclaration ::= "type" Id "=" Type ";"
// <: is not real TS syntax, defined here to mirror the checkTypeRelatedTo behavior (aka isAssignable)
// /*!\url{https://raw.githubusercontent.com/microsoft/TypeScript/refs/heads/main/src/compiler/checker.ts#:~:text=is%20truthy.%0A%20%20%20%20%20*/%0A%20%20%20%20function-,checkTypeRelatedTo,-(%0A%20%20%20%20%20%20%20%20source%3A%20Type%2C%0A%20%20%20%20%20%20%20%20target}!*/
IsAssignableTo ::= Type "<:" Type ";"
Stmt ::= TypeAliasDeclaration
       | IsAssignableTo
Pgm ::= List{Stmt, ""}
\end{lstlisting}

\section{TypeScript Type Checking Environment Configuration (WIP)}

$\Gamma := \mathtt{Identifier}\mapsto \mathtt{Type}$

\section{TypeScript Type Checking Rules (WIP)}

\begin{prooftree}
\AxiomC{}
\UnaryInfC{$\Gamma, \texttt{Id} \mapsto \texttt{T} | \texttt{Id} \Downarrow \texttt{T}$}{LookUp}
\end{prooftree}

\begin{prooftree}
\AxiomC{}
\UnaryInfC{$\Gamma | \texttt{ type Id = T ;} \Downarrow \Gamma, \texttt{Id} \mapsto \texttt{T}$}{Definition}
\end{prooftree}

\subsection*{Program sequencing and result accumulation}

Currently, we evaluate a program as a list of statements, accumulating a list $R$ of results produced by $\texttt{<:}$ checks.

\subsection*{Core assignability ($<:$) rules}

We write the judgment $S \mathrel{<:} T$ for $\mathrm{isAssignable}(S, T)$.

\begin{multicols}{2}
\begin{prooftree}
\AxiomC{}
\UnaryInfC{$T \mathrel{<:} T$}{Sub-Refl}
\end{prooftree}

\begin{prooftree}
\AxiomC{}
\UnaryInfC{$\texttt{any} \mathrel{<:} T$}{Any-Top-L}
\end{prooftree}

\begin{prooftree}
\AxiomC{}
\UnaryInfC{$T \mathrel{<:} \texttt{any}$}{Any-Top-R}
\end{prooftree}

\begin{prooftree}
\AxiomC{}
\UnaryInfC{$S \mathrel{<:} \texttt{unknown}$}{Unknown-Top}
\end{prooftree}

\begin{prooftree}
\AxiomC{}
\UnaryInfC{$\texttt{unknown} \mathrel{<:} \texttt{any}$}{Unknown-Any}
\end{prooftree}

\begin{prooftree}
\AxiomC{}
\UnaryInfC{$\texttt{unknown} \mathrel{<:} \texttt{unknown}$}{Unknown-Refl}
\end{prooftree}


\begin{prooftree}
\AxiomC{}
\UnaryInfC{$\texttt{never} \mathrel{<:} T$}{Never-Bot}
\end{prooftree}

% Primitive widening rules (string-like, number-like, etc.)

\begin{prooftree}
\AxiomC{}
\UnaryInfC{$\texttt{string} \mathrel{<:} \texttt{string}$}{String-Refl}
\end{prooftree}

\begin{prooftree}
\AxiomC{}
\UnaryInfC{$\texttt{number} \mathrel{<:} \texttt{number}$}{Number-Refl}
\end{prooftree}

\begin{prooftree}
\AxiomC{}
\UnaryInfC{$\texttt{bigint} \mathrel{<:} \texttt{bigint}$}{BigInt-Refl}
\end{prooftree}

\begin{prooftree}
\AxiomC{}
\UnaryInfC{$\texttt{boolean} \mathrel{<:} \texttt{boolean}$}{Boolean-Refl}
\end{prooftree}

\begin{prooftree}
\AxiomC{}
\UnaryInfC{$\texttt{symbol} \mathrel{<:} \texttt{symbol}$}{Symbol-Refl}
\end{prooftree}

\begin{prooftree}
\AxiomC{$S \mathrel{<:} T_1$}
\UnaryInfC{$S \mathrel{<:} T_1 \mid T_2$}{Union-R1}
\end{prooftree}

\begin{prooftree}
\AxiomC{$S \mathrel{<:} T_2$}
\UnaryInfC{$S \mathrel{<:} T_1 \mid T_2$}{Union-R2}
\end{prooftree}

\begin{prooftree}
\AxiomC{$S_1 \mathrel{<:} T$}
\AxiomC{$S_2 \mathrel{<:} T$}
\BinaryInfC{$S_1 \mid S_2 \mathrel{<:} T$}{Union-L}
\end{prooftree}
\end{multicols}

\paragraph{Non-derivable judgments.}
There are no rules for $\texttt{unknown} \mathrel{<:} T$ when $T \not\in \{\texttt{any},\texttt{unknown}\}$, nor for $T \mathrel{<:} \texttt{never}$ unless $T=\texttt{never}$; thus those judgments are not derivable.

% TODO: Extend the formalization to cover additional cases handled by
% `isSimpleTypeRelatedTo` in the TypeScript checker:
%   * enum types and enum-literal types
%   * literal-to-base conversions (e.g., "foo" <: string, 1 <: number)
%   * null/undefined special-case rules (including non-strict null checks
%     and unions that behave like "unknown-like")
% These are currently omitted from the core rules for simplicity.

\paragraph{Intersection types.}

For intersection types, we adopt the usual greatest-lower-bound reading: a value of type
$S_1 \,\&\, S_2$ satisfies both $S_1$ and $S_2$. At the level of assignability, we capture
this with the following rule:

\begin{prooftree}
\AxiomC{$S_1 \,\&\, S_2$}
\AxiomC{$S_1 \mathrel{<:} T$}
\AxiomC{$S_2 \mathrel{<:} T$}
\TrinaryInfC{$S_1 \,\&\, S_2 \mathrel{<:} T$}{Inter-L}
\end{prooftree}

That is, an intersection type is assignable to a target type $T$ exactly when each
constituent is assignable to $T$. Note that we intentionally do \emph{not} add a converse
rule of the form $S \mathrel{<:} T_1 \,\&\, T_2$, since TypeScript's implementation of
intersections (via normalization and structural checks) does not behave as a simple
greatest lower bound in all cases.

\subsection*{Object structural subtyping}


We write object types in the form
\[
  \{ \boldsymbol{P} \}
\]
where each property descriptor $P$ is of one of the following forms:
\[
  x : T \quad\mid\quad x? : T \quad\mid\quad \texttt{readonly}\ x : T \quad\mid\quad \texttt{readonly}\ x? : T.
\]

We use the following meta-level predicates:
\begin{itemize}
  \item $\mathrm{name}(P)$ is the identifier of a property descriptor.
  \item $\mathrm{type}(P)$ is its type.
  \item $\mathrm{readonly}(P)$ is true iff $P$ is declared \texttt{readonly}.
  \item $\mathrm{required}(P)$ is true iff $P$ is not marked with ``\texttt{?}''.
\end{itemize}

\paragraph{Structured object subtyping (ignoring functions and intersections).}

Let $S = \{ \boldsymbol{P_s} \}$ and $T = \{ \boldsymbol{P_t} \}$.
We say that $S \mathrel{<:} T$ holds when all of the following conditions are satisfied:

\begin{itemize}
  \item For every property descriptor $P_t \in \boldsymbol{P_t}$ with name $x = \mathrm{name}(P_t)$, there exists \emph{some}
        $P_s \in \boldsymbol{P_s}$ with $\mathrm{name}(P_s) = x$ such that:
        \begin{itemize}
          \item If $\mathrm{required}(P_t)$ then $\mathrm{required}(P_s)$.
          \item If $\mathrm{readonly}(P_t)$ then $\mathrm{readonly}(P_s)$.
          \item $\mathrm{type}(P_s) \mathrel{<:} \mathrm{type}(P_t)$.
        \end{itemize}
  \item No further condition is imposed on properties of $S$ whose names do not appear in $T$ (extra properties on $S$ are allowed).
\end{itemize}

This can be encoded by the following rule schema:

\begin{prooftree}
\AxiomC{$\forall P_t \in \boldsymbol{P_t}.\ \exists P_s \in \boldsymbol{P_s}.\
  \begin{array}{l}
    \mathrm{name}(P_s) = \mathrm{name}(P_t)\ \wedge\\
    \bigl(\mathrm{required}(P_t) \Rightarrow \mathrm{required}(P_s)\bigr)\ \wedge\\
    \bigl(\mathrm{readonly}(P_t) \Rightarrow \mathrm{readonly}(P_s)\bigr)\ \wedge\\
    \bigl(\mathrm{type}(P_s) \mathrel{<:} \mathrm{type}(P_t)\bigr)
  \end{array}$}
\UnaryInfC{$\{ \boldsymbol{P_s} \} \mathrel{<:} \{ \boldsymbol{P_t} \}$}
\end{prooftree}


\end{document}
